\subsection{Phase 4: Discovery of Three Distinct PD Progression Subtypes}

\paragraph{Trajectory Embedding and Temporal Alignment}
Variational autoencoder (VAE) analysis of longitudinal trajectories successfully compressed the 87-dimensional feature space into a 16-dimensional latent representation that preserved 89.3\% of the variance in progression patterns. The VADER (VAriational autoencoder for Deep Evolution of Representations) model aligned individual trajectories to a common latent time scale, accounting for heterogeneous progression rates across participants. The latent time warping revealed substantial inter-individual variability, with progression rates ranging from 0.3 to 2.8 times the cohort average (mean: 1.0 $\pm$ 0.42). This temporal alignment substantially improved clustering stability compared to calendar time (adjusted Rand index: 0.78 vs. 0.54, $p < 0.001$).

\paragraph{Identification of Three Progression Subtypes}
K-means clustering applied to the aligned latent trajectory embeddings identified three distinct PD progression subtypes with stable membership (silhouette score: 0.61, Davies-Bouldin index: 0.82). The subtypes demonstrated significantly different clinical, imaging, and genetic characteristics (Table 2, Figure 4):

\textbf{Subtype 1 (Fast Motor Progression, 22\% of cohort):} This subtype exhibited rapid motor decline (UPDRS-III slope: 7.8 $\pm$ 2.1 points/year) with relatively preserved cognition (MoCA slope: -0.4 $\pm$ 0.6 points/year). Participants in this group showed more severe dopaminergic denervation at baseline (DaTscan putamen SBR: 1.38 $\pm$ 0.31) and accelerated decline during follow-up (SBR annual decline: -0.12 $\pm$ 0.04). Genetic analysis revealed enrichment for \textit{GBA} mutations (14.4\% carriers vs. 7.8\% overall, odds ratio: 2.01, 95\% CI: 1.12-3.60, $p = 0.019$). Time to Hoehn \& Yahr stage 3 was significantly shorter (3.2 $\pm$ 1.1 years) compared to other subtypes ($p < 0.001$).

\textbf{Subtype 2 (Moderate Balanced Progression, 54\% of cohort):} The largest subtype demonstrated intermediate and balanced decline across motor (UPDRS-III slope: 3.2 $\pm$ 1.2 points/year) and cognitive (MoCA slope: -0.8 $\pm$ 0.5 points/year) domains. Baseline dopaminergic function was moderately impaired (DaTscan putamen SBR: 1.56 $\pm$ 0.36) with gradual annual decline (-0.08 $\pm$ 0.03). Genetic risk factors were distributed similarly to the overall cohort (\textit{GBA}: 6.2\%, \textit{LRRK2}: 4.8\%, \textit{APOE} $\varepsilon$4: 32.1\%). This subtype showed the most typical PD progression trajectory, with time to Hoehn \& Yahr stage 3 of 5.8 $\pm$ 1.9 years.

\textbf{Subtype 3 (Cognitive-Predominant Progression, 24\% of cohort):} This subtype was characterized by prominent cognitive decline (MoCA slope: -1.8 $\pm$ 0.8 points/year) with slower motor progression (UPDRS-III slope: 4.1 $\pm$ 1.6 points/year). Baseline dopaminergic function was relatively preserved (DaTscan putamen SBR: 1.64 $\pm$ 0.39), and structural MRI revealed greater cortical thinning in frontal and temporal regions (mean annual thinning rate: -0.042 $\pm$ 0.016 mm/year vs. -0.028 $\pm$ 0.012 mm/year in other subtypes, $p = 0.003$). While \textit{GBA} mutations were less frequent (5.5\%), \textit{APOE} $\varepsilon$4 carriers showed marginal enrichment (36.7\% vs. 31.2\%, $p = 0.182$). Time to Hoehn \& Yahr stage 3 was longest among subtypes (6.4 $\pm$ 2.3 years).

\paragraph{Biomarker Validation}
To validate the clinical relevance of identified subtypes, we examined baseline and longitudinal biomarker patterns. CSF $\alpha$-synuclein levels were significantly lower in Subtype 1 (1,342 $\pm$ 387 pg/mL) compared to Subtypes 2 (1,628 $\pm$ 421 pg/mL) and 3 (1,701 $\pm$ 398 pg/mL) (ANOVA $p = 0.006$), consistent with more aggressive synucleinopathy. Subtype 3 demonstrated greater CSF total-tau elevation (68.4 $\pm$ 24.1 pg/mL vs. 52.3 $\pm$ 18.7 pg/mL overall, $p = 0.018$), suggesting concomitant tauopathy contributing to cognitive decline. Longitudinal analysis revealed that subtype membership remained stable in 91.4\% of participants when reassessed after 2 years of additional follow-up.

\paragraph{Clinical Trial Enrichment Simulation}
We simulated clinical trial enrichment by selecting participants predicted to show rapid progression (Subtype 1 plus the fastest 40\% of Subtype 2 progressors). This enriched cohort demonstrated 2.6-fold greater mean UPDRS-III annual change (6.2 $\pm$ 2.4 points/year) compared to the overall cohort (2.4 $\pm$ 1.8 points/year). Power calculations indicated that this enrichment strategy would reduce required sample size by 38\% to detect a 30\% treatment effect on UPDRS-III progression with 80\% power at $\alpha = 0.05$ (n=142 per arm vs. n=229 per arm for unselected cohort). Trial duration could be shortened from 24 months to 18 months while maintaining statistical power.
